\chapterdivider{1}{Orientation and application map}
  {What can JaxCont compute, and where should a newcomer begin?}
  {Featured application: one model viewed as branches, parameter curves, cycles, phase geometry, and phase sensitivity}

\begin{frame}{How to read ``v0.4.0''}
\small
\begin{columns}[T,onlytextwidth]
\column{.30\textwidth}
\releasedbadge
\column{.67\textwidth}
Continuation, periodic orbits, Floquet events, Hopf classification, direct
codim-2 solvers, phase planes, PRC/dPRC, and the validation suite.
\end{columns}

\vspace{4mm}
\begin{columns}[T,onlytextwidth]
\column{.30\textwidth}
\releasedbadge
\column{.67\textwidth}
Fold/Hopf two-parameter continuation, their codimension-two events,
\texttt{plot\_two\_parameter\_diagram}, and Examples 15--20.
\end{columns}

\vspace{4mm}
\begin{columns}[T,onlytextwidth]
\column{.30\textwidth}
\futurebadge
\column{.67\textwidth}
Branch switching, general BVP and connecting-orbit workflows, and
periodic-orbit codimension-two curve continuation.
\end{columns}

\vfill
\begin{alertblock}{Release scope}
The green badges identify capabilities included in v0.4.0. Red badges mark
workflows that remain outside the supported surface.
\end{alertblock}
\end{frame}

\begin{frame}{Start from the scientific question}
\centering
\begin{tikzpicture}[node distance=1.6mm and 8mm,>=Latex,
 question/.style={rounded corners,draw=JaxBlue,thick,fill=SoftBlue,
   minimum width=66mm,minimum height=7mm,align=left,font=\small},
 method/.style={rounded corners,draw=JaxTeal,thick,fill=JaxTeal!9,
   minimum width=56mm,minimum height=7mm,align=center,font=\small}]
\node[question] (q1) {Where do steady states exist?};
\node[method,right=of q1] (m1) {equilibrium continuation};
\node[question,below=of q1] (q2) {Where do oscillations exist?};
\node[method,right=of q2] (m2) {periodic-orbit continuation};
\node[question,below=of q2] (q3) {Where does stability change?};
\node[method,right=of q3] (m3) {eigenvalues, Floquet multipliers, events};
\node[question,below=of q3] (q4) {What kind of local organizing point is this?};
\node[method,right=of q4] (m4) {Hopf $l_1$ and direct codim-2 refinement};
\node[question,below=of q4] (q5) {How does an event move in a parameter plane?};
\node[method,right=of q5] (m5) {fold/Hopf curve continuation};
\node[question,below=of q5] (q6) {How does an oscillator respond to perturbations?};
\node[method,right=of q6] (m6) {iPRC and dPRC};
\foreach \i in {1,...,6}{\draw[->,thick] (q\i)--(m\i);}
\end{tikzpicture}

\vfill
\begin{block}{Application map}
Choose the object or change you want to understand; that question selects the
continuation problem, diagnostic, or sensitivity calculation.
\end{block}
\end{frame}

\begin{frame}{Choose a route through the presentation}
\begin{columns}[T,onlytextwidth]
\column{.32\textwidth}
\begin{block}{Overview route}
\centering
\small Chapters

\Large 1, 2

\normalsize selected 3, 5, 8--9

\Large 12, 13

\vspace{2mm}
\small Intuition, representative outputs, workflow, and boundaries.
\end{block}

\column{.32\textwidth}
\begin{block}{Methods route}
\centering
\small Chapters

\Large 2--8, 11

\vspace{8.4mm}
\small Algorithms, stability, local refinements, and validation.
\end{block}

\column{.32\textwidth}
\begin{block}{Complete route}
\centering
\small Chapters

\Large 1--13

\vspace{8.4mm}
\small The full narrative, including applications and guided workflows.
\end{block}
\end{columns}

\vfill
\centering
\textcolor{JaxTeal}{Chapter 13 remains available as the release-integrity, scope, and glossary reference.}
\end{frame}

\begin{frame}{What you should be able to explain afterwards}
\footnotesize
\begin{columns}[T,onlytextwidth]
\column{.49\textwidth}
\begin{enumerate}
  \item Distinguish continuation from simulation and isolated root solving.
  \item Explain fold failure and pseudo-arclength passage.
  \item Build equilibrium and periodic-orbit problems with the public API.
  \item Interpret eigenvalues, Floquet multipliers, and codim-1 events.
  \item Explain Hopf refinement and the first Lyapunov coefficient.
\end{enumerate}

\column{.49\textwidth}
\begin{enumerate}
  \setcounter{enumi}{5}
  \item State the purpose and limits of CP, BT, GH, ZH, and HH solvers.
  \item Trace fold/Hopf curves and interpret codimension-two curve events.
  \item Distinguish parameter-plane diagrams, branch plots, and 2D phase planes.
  \item Explain iPRCs and their parameter derivatives.
  \item Run examples, inspect outputs, and choose a validation strategy.
  \item State the boundaries: no branch switching or periodic codim-2 curves.
\end{enumerate}
\end{columns}

\vfill
\begin{block}{The through-line}
Ask a scientific question, choose the numerical object, then validate the
result before interpreting it.
\end{block}
\end{frame}
